%%%%%%%%%%%%%%%%%%%%%%%%%%%%%%%%%%%%%%%%%%%%%%%%%%%%%%%%%%%%%%%%%%%%%%
% Lớp 9 | Chương V. Đường tròn | Độ dài của cung tròn. Diện tích hình quạt tròn và hình vành khuyên
% Dạng 9H5?15-2: Diện tích hình tròn, hình quạt tròn
%%%%%%%%%%%%%%%%%%%%%%%%%%%%%%%%%%%%%%%%%%%%%%%%%%%%%%%%%%%%%%%%%%%%%%


% nguồn: dan_tu_AI (nhập 14/11/2026)
\begin{ex}%[9H5N15-2]%[ma:31886191]
	Diện tích của hình tròn có bán kính $3\text{ cm}$ là
	\choice
	{\True $9\pi\text{ cm}^2$}
	{$6\pi\text{ cm}^2$}
	{$3\pi\text{ cm}^2$}
	{$36\pi\text{ cm}^2$}
	\loigiai{
		$S=\pi R^2=\pi\cdot 3^2=9\pi\text{ (cm}^2)$.
	}
\end{ex}

% nguồn: dan_tu_AI (nhập 14/11/2026)
\begin{ex}%[9H5H15-2]%[ma:b14474cd]
	Diện tích của hình quạt tròn bán kính $6\text{ cm}$, ứng với cung $60^\circ$ là
	\choice
	{$2\pi\text{ cm}^2$}
	{\True $6\pi\text{ cm}^2$}
	{$12\pi\text{ cm}^2$}
	{$3\pi\text{ cm}^2$}
	\loigiai{
		\[ S=\dfrac{\pi R^2n}{360}=\dfrac{\pi\cdot 6^2\cdot 60}{360}=6\pi\text{ (cm}^2). \]
	}
\end{ex}

% nguồn: dan_tu_AI (nhập 14/11/2026)
\begin{ex}%[9H5H15-2]%[ma:2f0a8dd9]
	Cho hình tròn có bán kính $4\text{ cm}$. Các khẳng định sau đúng hay sai?
	\choiceTF
	{\True Diện tích hình tròn bằng $16\pi\text{ cm}^2$}
	{Diện tích hình quạt tròn ứng với cung $45^\circ$ bằng $4\pi\text{ cm}^2$}
	{Hình quạt tròn có diện tích $4\pi\text{ cm}^2$ thì ứng với cung $45^\circ$}
	{Hình quạt tròn có độ dài cung $2\pi\text{ cm}$ thì có diện tích bằng $8\pi\text{ cm}^2$}
	\loigiai{
		\begin{itemchoice}
			\itemch Đúng. $S=\pi\cdot 4^2=16\pi\text{ (cm}^2)$.
			\itemch Sai. $S_q=\dfrac{\pi\cdot 4^2\cdot 45}{360}=2\pi\text{ (cm}^2)$.
			\itemch Sai. $\dfrac{16\pi n}{360}=4\pi\Leftrightarrow n=90$, nên hình quạt ứng với cung $90^\circ$.
			\itemch Sai. $S_q=\dfrac{lR}{2}=\dfrac{2\pi\cdot 4}{2}=4\pi\text{ (cm}^2)$.
		\end{itemchoice}
	}
\end{ex}
